\section{AI 安全与国际监管}

\subsection{两大核心风险}

Hassabis 将 AI 安全风险分为两个维度：

\begin{enumerate}
    \item \textbf{恶意使用（Misuse）}：坏人利用 AI 工具造成伤害——这是 Hassabis 认为的\textbf{首要安全担忧}（primary concern）
    \item \textbf{技术对齐（Technical Alignment）}：确保 AI 系统的行为与人类意图一致
\end{enumerate}

\begin{importantbox}{恶意使用是第一优先级}
Hassabis 将恶意使用——而非抽象的``AI 失控''——列为首要安全关切。这包括：
\begin{itemize}
    \item 使用 AI 进行大规模网络攻击
    \item 利用 AI 辅助生物武器设计
    \item 用 AI 生成深度伪造内容进行社会工程攻击
    \item 国家和非国家行为体利用 AI 增强破坏能力
\end{itemize}
\textit{关键洞察：最紧迫的威胁往往不是 AI 自己``变坏''，而是人类主动将 AI 用于恶意目的。}
\end{importantbox}

\subsection{``Kite Mark'' 认证与欺骗检测}

Hassabis 提出了一套务实的安全验证框架：

\begin{itemize}
    \item \textbf{欺骗行为基准测试}（benchmarks for deception）：系统性地检测 AI 模型是否具备欺骗能力或倾向
    \item \textbf{``Kite Mark'' 认证}：类似产品安全认证标志——通过独立测试的模型获得安全认证标记
    \item \textbf{AI Safety Institutes 的验证角色}：支持各国 AI 安全研究所承担独立的第三方验证职能
\end{itemize}

\begin{knowledgebox}{什么是 ``Kite Mark''？}
Kite Mark 是英国标准协会（BSI）颁发的产品安全认证标志，类似于中国的 CCC 认证或欧盟的 CE 标志。Hassabis 用这个类比来说明：AI 模型也需要类似的、由独立机构颁发的安全认证——不是开发者自己说``我的模型是安全的''，而是经过标准化测试后由第三方背书。
\end{knowledgebox}

\subsection{国际监管体系的设想}

Hassabis 的首要政策诉求是建立一个\textbf{类似国际原子能机构（IAEA）的 AI 国际监管机构}。

\begin{importantbox}{Hassabis 的政策首选：AI 版 IAEA}
\textit{Hassabis 最希望看到的政策行动是：建立一个类似 IAEA 的国际机构，专门负责 AI 安全的全球协调与验证。}

这个类比的深层含义：
\begin{itemize}
    \item \textbf{核武器的教训}：核技术是人类历史上第一个需要国际协调的``存在性风险''技术。IAEA 的成立花了十多年，但最终证明了其必要性。
    \item \textbf{生物技术的教训}：《生物武器公约》（BWC）虽然有缺陷，但提供了国际规范的基本框架。
    \item \textbf{AI 的特殊性}：AI 比核技术更容易扩散（不需要铀矿），因此国际协调更加紧迫。
\end{itemize}
\end{importantbox}

\begin{warningbox}{为什么自律是不够的}
一个常见的行业论调是``让我们自己来监管自己就好了''。Hassabis 明确反对这种观点：

AI 安全不能仅依赖企业自律——需要有强制力的国际标准和独立验证机制。原因很简单：在竞争压力下，没有外部约束的自律承诺是脆弱的。就像没有哪个核电站会说``我们不需要 IAEA 的检查，我们自己就能保证安全''。
\end{warningbox}

\subsection{本章小结}

AI 安全的首要威胁是恶意使用，需要通过欺骗检测基准和``Kite Mark''式认证来建立信任。Hassabis 最核心的政策诉求是建立 AI 版 IAEA——一个具有国际约束力的安全监管与验证机构。

%% ==================== Section 8 ====================
